% RESTORED 2026-09-22 (v3.63): pasted back into chapters/06_results.tex on the pillars_hg fallback
% after Gen15 U17 (pillars_xl) was abandoned. This copy is history; the chapter is the text of record.
% =============================================================================
%  WITHHELD FROM THE DRAFT -- 2026-09-18 (v3.40)
%  UAV-pillars results, cut out of chapters/06_results.tex verbatim.
%
%  WHY: the scene does not test what the chapter uses it for. Its constraint set
%  is the one the demonstration generator was built to satisfy -- the channels sit
%  at |y| = 1.11 m precisely because 0.6 + 0.12 + rotor reach + safety = 1.11 --
%  so the plans are already feasible and the projection has nothing to repair.
%  See data_status/PENDING_20260918_pillars_geometry_redesign.md.
%
%  RESTORE: paste back between \subsubsection{UAV-corridor} ... and
%  \subsubsection{UAV-s-curve}, once the enlarged-constraint evaluations exist,
%  with every number recomputed for the new geometry.
% =============================================================================

\subsubsection{UAV-pillars}\label{sec:res:uav:pillars}

UAV-pillars is the scene on which the projection has the least to repair, and reading it requires saying
that first. Executed with no projection at all, instantaneous-velocity matching passes the goal
collision-free on all ten flights and analytic average-velocity matching on nine
(\autoref{tab:uav-pillars}, first column). The plans are already feasible. What the projected rows then
measure is not whether a method can make an infeasible plan feasible --- the question UAV-corridor
answers --- but how much of an already feasible plan each method preserves. Consistency-interpolated
average-velocity matching and the diffusion baseline are the two that do need repair: neither finishes a
single clean flight unprojected, and the baseline crosses the line on every flight while violating the
pillar clearance on $132$ control steps per flight.

The evaluation covers two blocks of configurations, and they behave differently enough that averaging
across them hides the result. The per-step block is seven configurations --- the plan without projection,
and per-step projection with each of the three selection rules on plain and on tightened constraints ---
and is evaluated at every budget that exists for a model. The endpoint block is four configurations, a
single candidate and each of the three rules, and exists only at $\nfe=5$.
\autoref{tab:uav-pillars} reports the first and \autoref{tab:uav-pillars-endpoint} the second.

\begin{table}[htpb]
  \caption[UAV-pillars: per-step projection at every evaluated budget]{UAV-pillars, the per-step block:
  success with constraint satisfaction --- the goal passed on a collision-free flight --- ten flights per
  cell, full constraint set, geometry \texttt{pillars\_hg}. Every budget at which a model was evaluated on
  this geometry is listed. $r$, $c$, $t$: random selection, cumulative projection cost, temporal
  consistency. \emph{tight.}: the tightened constraint set. The mean is over the seven configurations of
  the row, which are the same seven in every row. Consistency-interpolated average-velocity matching with
  $\alpha_{\mathrm{end}}=0.2$. Temporal U-Net, $3.96$--$3.97$\,M parameters; training seed 6.
  ---: not evaluated.}
  \label{tab:uav-pillars}
  \centering
  \footnotesize
  \setlength{\tabcolsep}{4pt}
  \begin{tabularx}{\linewidth}{@{}L r rrrr rrr r@{}}
    \toprule
    & & no & \multicolumn{3}{c}{per-step} & \multicolumn{3}{c}{per-step tight.} & \\
    \cmidrule(lr){4-6}\cmidrule(lr){7-9}
    Model & $\nfe$ & proj. & $r$ & $c$ & $t$ & $r$ & $c$ & $t$ & mean \\
    \midrule
    MeanFM    & 2  & 0.30 & 0.70 & 0.40 & 0.40 & 0.40 & 0.40 & 0.40 & 0.429 \\
    MeanFM    & 5  & 0.90 & 0.20 & 0.80 & 0.40 & 0.10 & 0.90 & 0.80 & 0.586 \\
    \midrule
    CI-MeanFM & 1  & 0.00 & 0.00 & 0.00 & 0.20 & 0.20 & 0.00 & 0.70 & 0.157 \\
    CI-MeanFM & 2  & 0.20 & 0.20 & 0.20 & 0.50 & 0.30 & 0.30 & 0.40 & 0.300 \\
    CI-MeanFM & 5  & 0.00 & 0.10 & 0.90 & 0.80 & 0.00 & 0.00 & 0.90 & 0.386 \\
    \midrule
    FM        & 2  & 0.30 & 0.50 & 0.20 & 0.30 & 0.40 & 0.50 & 0.20 & 0.343 \\
    FM        & 5  & 1.00 & 0.00 & 1.00 & 0.80 & 0.00 & 0.90 & 0.40 & 0.586 \\
    \midrule
    Diffusion & 20 & 0.00 & ---  & 0.00 & 0.00 & ---  & 0.10 & 0.30 & 0.080 \\
    \bottomrule
  \end{tabularx}
\end{table}
\dataref{analysis\_results\_checkpoint/15-09/batch\_uav\_20260915\_100816, geometry
\texttt{pillars\_hg\_bounds+dynamics+geo\_bounds+obstacles}, evaluation tags \texttt{u7hg} and
\texttt{EPlatest\_u7hg}, mask \texttt{all}; recomputed cell by cell by
\texttt{DA\_in\_Paper/analysis/pillars\_grid.py}, which resolves the model of every candidate through
\texttt{candidate\_axes.csv} (\texttt{Full\_Path}) and excludes the \texttt{geo\_free} variants}
\hole{MeanFM and FM at $\nfe=1$ on this geometry. Both exist at $\nfe=1$ only on the earlier
\texttt{pillars} geometry, which is a different constraint set and cannot be put in this table.
CI-MeanFM is the only model with an $\nfe=1$ row here, so the first budget carries no comparison.}
\guard{The diffusion baseline was not evaluated with random selection, plain or tightened; its mean is
over the five configurations it has.}

\begin{table}[htpb]
  \caption[UAV-pillars: endpoint projection]{UAV-pillars, the endpoint block: the same measure under
  endpoint projection, which exists at $\nfe=5$ only. \emph{single}: one candidate; $r$, $c$, $t$: the
  three selection rules over four candidates. The diffusion baseline has no endpoint row on this scene.}
  \label{tab:uav-pillars-endpoint}
  \centering
  \footnotesize
  \setlength{\tabcolsep}{6pt}
  \begin{tabularx}{\linewidth}{@{}L r rrrr r@{}}
    \toprule
    Model & $\nfe$ & single & $r$ & $c$ & $t$ & mean \\
    \midrule
    MeanFM    & 5  & 0.50 & 0.90 & 0.70 & 0.70 & 0.700 \\
    CI-MeanFM & 5  & 0.20 & 0.40 & 0.40 & 0.30 & 0.325 \\
    FM        & 5  & \mainconfig{\selectedmark\ 1.00} & 1.00 & 0.90 & 0.90 & \textbf{0.950} \\
    Diffusion & 20 & ---  & ---  & ---  & ---  & --- \\
    \bottomrule
  \end{tabularx}
\end{table}
\dataref{same batch and geometry as \autoref{tab:uav-pillars}, variants \texttt{hardflow\_sls}
and \texttt{hardflow\_sls-\{r,c,t\}}}

\paragraph{What the two blocks say.} On the seven per-step configurations at $\nfe=5$, analytic
average-velocity matching and instantaneous-velocity matching are level: both average $0.586$, both reach
$0.80$ or more in four of the seven, and each is ahead of the other in three of the seven cells.
Consistency-interpolated average-velocity matching averages $0.386$ and the diffusion baseline $0.080$.
The separation between the two leading models appears only in the endpoint block, where
instantaneous-velocity matching averages $0.950$ against $0.700$ --- and that block sits on top of the
plan that was already collision-free on all ten flights before any projection ran. At $\nfe=2$, where
only the per-step block exists, the order is the other way round: analytic average-velocity matching
$0.429$, instantaneous-velocity matching $0.343$, consistency-interpolated $0.300$.

So this scene supports a narrower claim than a single mean would suggest. It does not show that
instantaneous-velocity matching is the better constrained planner here; it shows that on a scene where
its unprojected plan is already feasible, endpoint projection at $\nfe=5$ leaves that plan intact while
per-step projection does not, and that at the smaller budget the ordering reverses. The one ordering that
holds in every cell of both blocks is against the diffusion baseline, which reaches $0.30$ at best and
finishes no flight within the constraints in three of its five configurations.

\autoref{tab:uav-pillars-best} places each model at its best configuration \emph{within each block and at
each budget it was evaluated at}, which is the only way the baseline enters the comparison on equal terms:
it has no endpoint row on this scene, so a single best-of-eleven column would have compared three models
on eleven configurations against a baseline that had five.

Read that way the scene gives no single winner. At $\nfe=5$, instantaneous-velocity matching reaches
$1.00$ in both blocks, analytic average-velocity matching and consistency-interpolated average-velocity
matching reach $0.90$ with per-step projection, and analytic average-velocity matching also reaches $0.90$
with endpoint projection at a tenth of the time per control step. At $\nfe=2$, where only per-step
projection was run, the order is reversed: analytic average-velocity matching $0.70$, instantaneous-velocity
matching and consistency-interpolated average-velocity matching $0.50$. Consistency-interpolated
average-velocity matching is the one model whose endpoint block is clearly worse than its per-step block,
$0.40$ against $0.90$. The baseline reaches $0.30$ and costs $4.4$\,s per control step.

\begin{table}[htpb]
  \caption[UAV-pillars: the best configuration of each model, per block and budget]{UAV-pillars: for every
  model and every budget evaluated on this geometry, the best configuration \emph{within} each projection
  block --- per-step (seven configurations) and endpoint (four). Comparing per block is what keeps the
  diffusion baseline in the comparison: it has no endpoint row here, and neither has any model at
  $\nfe\le2$. \emph{S\&C}: the goal passed on a collision-free flight; every row passed the goal on all ten
  flights, so the column is omitted. \emph{viol.}: violating control steps per flight; \emph{steps}: control
  steps, the episode ending after $634$. Ten flights per cell, training seed 6. A bold filled circle marks
  the baseline of record. No configuration is marked as selected: the scene does not support one
  (\autoref{sec:res:uav:conclusion}).}
  \label{tab:uav-pillars-best}
  \centering
  \footnotesize
  \setlength{\tabcolsep}{3pt}
  \begin{tabularx}{\linewidth}{@{}L r L rrr L rrr@{}}
    \toprule
    & & \multicolumn{4}{c}{best per-step configuration} & \multicolumn{4}{c}{best endpoint configuration} \\
    \cmidrule(lr){3-6}\cmidrule(lr){7-10}
    Model & $\nfe$ & Rule & S\&C & viol. & steps & Rule & S\&C & viol. & steps \\
    \midrule
    MeanFM    & 2  & $r$              & 0.70 & 54.7 & 491 & ---            & ---  & ---   & --- \\
    MeanFM    & 5  & tight., $c$      & 0.90 & 20.1 & 432 & $r$            & 0.90 & 21.4  & 442 \\
    \midrule
    CI-MeanFM & 1  & tight., $t$      & 0.70 & 53.7 & 592 & ---            & ---  & ---   & --- \\
    CI-MeanFM & 2  & $t$              & 0.50 & 91.3 & 528 & ---            & ---  & ---   & --- \\
    CI-MeanFM & 5  & tight., $t$      & 0.90 & 22.4 & 453 & $c$            & 0.40 & 113.8 & 616 \\
    \midrule
    FM        & 2  & $r$              & 0.50 & 97.7 & 538 & ---            & ---  & ---   & --- \\
    FM        & 5  & $c$              & 1.00 & 0.0  & 558 & single cand.   & 1.00 & 0.0   & 562 \\
    \midrule
    \mainconfig{\baselinemark\ Diffusion} & \mainconfig{20} & \mainconfig{tight., $t$} & \mainconfig{0.30} & \mainconfig{101.0} & \mainconfig{634} & \mainconfig{---} & \mainconfig{---} & \mainconfig{---} & \mainconfig{---} \\
    \bottomrule
  \end{tabularx}
\end{table}
\dataref{same batch, geometry and tags as \autoref{tab:uav-pillars}; best configuration per block and
budget by \texttt{DA\_in\_Paper/analysis/pillars\_grid.py::best\_per\_block}. Time per control step, not
shown above: MeanFM $60.3$ and $1594.0$\,ms per-step, $154.0$\,ms endpoint; CI-MeanFM $85.0$, $75.1$,
$1564.5$ and $253.4$\,ms; FM $78.4$ and $790.8$\,ms per-step, $78.0$\,ms endpoint; the baseline
$4424.5$\,ms}

The two rows that reach $0.90$ and $1.00$ at $\nfe=5$ have to be read flight by flight, because the rate
hides how they fly. Instantaneous-velocity matching with endpoint projection on a single candidate passes
the goal collision-free on all ten flights with no violating step --- but six of those ten run to the
$634$-step limit, crossing the line and never settling, and end $0.36$ to $0.39$\,m from the goal point.
Analytic average-velocity matching with endpoint projection and random selection does so on nine: eight of
its flights finish in $414$ to $435$ control steps and end $0.29$ to $0.30$\,m from the goal point, and one
flight loses the field, violates the clearance on $214$ steps and runs to the limit. That single flight is
the whole difference in the rate. So instantaneous-velocity matching is the model that never fails outright
here, and analytic average-velocity matching the one that flies the shorter, tighter path when it succeeds;
neither statement is a ranking of the two.

\autoref{fig:uav-pillars-paths} draws those four configurations, and it shows what the table cannot: the
models do not fail by stopping short, they fail by committing to a side. Instantaneous-velocity matching
sends all ten flights above the pillar field and keeps all ten clean; analytic average-velocity matching
and consistency-interpolated average-velocity matching send nine below and one above, and it is the one
above that clips a pillar. The diffusion baseline flies the upper lane like instantaneous-velocity
matching but is collision-free on three flights of ten, touching the pillars on the way.

\begin{figure}[htpb]
  \centering
  \includegraphics[width=\textwidth]{fig_uav_pillars_paths}
  \caption[UAV-pillars: flown paths at each model's configuration]{UAV-pillars seen from above, drawn like
  \autoref{fig:constraints-uav}: the six pillars inflated by the vehicle radius and tightened, with the
  executed path of each of the ten flights over them. Colour reports the constraints (green:
  collision-free; red: entered an inflated pillar) and the stroke and end mark report the goal (solid and
  filled: passed; dashed and hollow: not passed). The four panels are \emph{not} under the same
  projection: each model is drawn at one of the two configurations \autoref{tab:uav-pillars-best} reports as
  its best at $\nfe=5$ --- the endpoint one for analytic average-velocity matching and for
  instantaneous-velocity matching, the per-step one for consistency-interpolated average-velocity matching
  and for the baseline, which has no endpoint row --- named in the panel subtitle. The panel titles name the generative objectives: average-velocity matching
  is MeanFM, instantaneous-velocity matching is FM and consistency-interpolated is CI-MeanFM.}
  \label{fig:uav-pillars-paths}
\end{figure}
\dataref{DA\_in\_Paper/data/uav\_paths.json, extract/uav\_paths.py; four cells of
batch\_uav\_20260915\_100816, tag u7hg}

\paragraph{Step budget.} Every model improves with the budget on this scene, as the mean column of
\autoref{tab:uav-pillars} reads down each block. What the table does not carry is the price: over those
same seven configurations the time per control step rises from $69$--$79$\,ms at $\nfe=2$ to
$0.9$--$1.7$\,s at $\nfe=5$. Endpoint projection avoids that rise --- it is the block in which
instantaneous-velocity matching reaches $1.00$ at $78$\,ms per step --- which is the trade
\autoref{sec:res:uav:projection} takes up.
\dataref{same batch, geometry and tags as \autoref{tab:uav-pillars}; \texttt{avg\_time\_ms} over the
seven per-step configurations at each budget}

\paragraph{Consistency-interpolated average-velocity matching.} It is level with analytic average-velocity
matching on UAV-corridor and behind both other flow-based models on UAV-pillars: the lowest mean of the
three at every budget, and no collision-free flight at all without projection, where the other two are at
$0.90$ and $1.00$. Its ceiling with projection is $0.90$, reached with per-step projection on tightened
constraints. With D3IL-aligning (\autoref{sec:res:aligning:consistency}) and D3IL-avoiding
(\autoref{tab:state-models}), the consistency-interpolated objective is level with the analytic one where
the task is saturated and behind it where the task separates the models; it is ahead of it in no
environment.
\dataref{uav\_results.py, pillars K1/K2/K5; Campaign CLOSURE\_20260910 \S2, \S4; DA\_20260910 T1}

